% ============================================================================
% Section 7: Statistical Analysis
% ============================================================================
\section{Statistical Analysis of the Discovery Corpus}
\label{sec:statistics}

\subsection{Corpus Summary}

As of 2026-08-11 19:30 UTC+2, the \RAMA{} engine has processed
481 of 698 manuscript pages (68.9\% completion).
Table~\ref{tab:corpus_summary} summarizes the results.

\begin{table}[H]
\centering
\caption{Summary statistics of the \RAMA{} discovery run
(snapshot at 2026-08-11 19:42 UTC+2; run ongoing).}
\label{tab:corpus_summary}
\begin{tabular}{lrr}
\toprule
\textbf{Metric} & \textbf{Count} & \textbf{Percentage} \\
\midrule
Total pages processed & $\sim$540 & 77.4\% of corpus \\
Total discoveries in database & 1737 & --- \\
Lean~4 Verified (\textsc{verified}) & 597 & 34.4\% of total \\
Lean~4 Failed (\textsc{failed}) & 1140 & 65.6\% of total \\
Stable Unitary vacua & 468 & 78.4\% of verified \\
Unstable (Tachyonic) & 129 & 21.6\% of verified \\
\bottomrule
\end{tabular}
\end{table}

\begin{figure}[H]
\centering
\includegraphics[width=\textwidth]{figures/energy_distribution.pdf}
\caption{(a)~Distribution of RAMA energy $E$ across all 597 verified
candidates. (b)~Energy vs.\ fit error colored by stability
classification. (c)~Proportion of stable vs.\ tachyonic
candidates.}
\label{fig:energy_dist}
\end{figure}

\subsection{Energy Distribution}

\begin{proposition}[Energy Range]
Among the 477 verified discoveries, the RAMA energy
$E$ ranges from $E_{\min} = 1.1272$ (Discovery~$\alpha$)
to $E_{\max} \approx 8.15$.  The median energy is
approximately $E_{\mathrm{med}} \approx 3.5$.
\end{proposition}

The fit error $I$ (the normalized $\ell_2$ residual
against target coefficients) shows a strong inverse
correlation with population size and generation count:
\begin{itemize}
  \item 10-page dry run (pop=25, gen=5): median $I \approx 0.45$.
  \item Full run (pop=50, gen=15): median $I \approx 0.25$.
  \item Best candidate: $I = 0.0705$ (7.05\% residual).
\end{itemize}

\subsection{Stability Classification}

The thermodynamic stability classification is based on
the sign of the effective central charge:
\begin{equation}
  \text{Equilibrium} =
  \begin{cases}
    \text{Stable (Unitary)} & \text{if } \ceff > 0, \\
    \text{Unstable (Tachyonic)} & \text{if } \ceff < 0, \\
    \text{Marginal} & \text{if } \ceff = 0.
  \end{cases}
\end{equation}

\begin{proposition}[Stability Ratio]
Of the 477 verified candidates, 378 (79.2\%) have
$\ceff > 0$ and are classified as \emph{Stable (Unitary)}.
101 (21.2\%) have $\ceff < 0$ and are classified as
\emph{Unstable (Tachyonic)}.
\end{proposition}

\begin{remark}
The preponderance of stable configurations ($\sim$79\%)
is consistent with the hypothesis that Ramanujan's
manuscript entries preferentially encode physically
viable partition functions.  However, this ratio
may also reflect a selection bias in the evolutionary
search, which penalizes high-complexity candidates
that tend to have negative $\ceff$.  A controlled
experiment with uniform random sampling would be
needed to distinguish these explanations.
\end{remark}

\subsection{Lean~4 Verification Rate}

The Lean~4 verification success rate improved dramatically
between the initial run (pop=25, gen=5) and the full run
(pop=50, gen=15):
\begin{itemize}
  \item \textbf{Initial full-corpus run} (August~9--10):
        0/1141 verified (0\%).  All candidates used the
        legacy auto-formalization template with broken
        tactic syntax.
  \item \textbf{10-page dry run} (August~11, after
        remediation): 8/10 verified (80\%).
  \item \textbf{Full high-generation run} (August~11,
        ongoing): 477/1618 unique candidates verified
        (29.5\%).
\end{itemize}

The 29.5\% verification rate reflects the fact that many
$q$-expansions converge to the same exponent vector
(database deduplication on \texttt{(eta\_exponents, q\_shift)}),
and the auto-formalization template was fixed after the
initial failed run.
